% LUA-DAG Technical Whitepaper
\documentclass[11pt,a4paper,twoside]{report}

\usepackage[utf8]{inputenc}
\usepackage[T1]{fontenc}
\usepackage{lmodern}
\usepackage{microtype}
\usepackage[margin=1in]{geometry}
\usepackage{amsmath,amssymb,amsthm}
\usepackage{booktabs}
\usepackage{longtable}
\usepackage{array}
\usepackage{multirow}
\usepackage{enumitem}
\usepackage{hyperref}
\usepackage{xcolor}
\usepackage{listings}
\usepackage{caption}
\usepackage{subcaption}
\usepackage{footnotehyper}

\hypersetup{
  colorlinks=true,
  linkcolor=blue!60!black,
  citecolor=blue!60!black,
  urlcolor=blue!70!black,
  pdftitle={LUA-DAG: Modular DA and Accountable Finality on a DAG},
  pdfauthor={LUA-DAG Contributors}
}

\lstset{
  basicstyle=\ttfamily\small,
  breaklines=true,
  frame=single,
  columns=fixed,
  keepspaces=true,
  showstringspaces=false
}

\theoremstyle{definition}
\newtheorem{definition}{Definition}[chapter]
\newtheorem{remark}{Remark}[chapter]

\title{LUA-DAG:\\Modular Data Availability and Accountable Finality\\on a Directed Acyclic Graph}
\author{}
\date{}

\begin{document}

% Avoid duplicate PDF page-1 anchor from \maketitle + abstract (hyperref)
\hypersetup{pageanchor=false}
\maketitle
\hypersetup{pageanchor=true}

\begin{abstract}
LUA-DAG is a \emph{consensus and data availability} protocol intended for rollups and application-specific chains. The design deliberately omits a general-purpose execution layer and does not aim to compete head-on with monolithic Layer~1 systems such as Sui, Aptos, or Monad. Instead, it targets the modular data availability (DA) plus shared finality segment, with Celestia, Avail, and EigenDA as natural points of comparison.

The architecture concentrates on two responsibilities: (i)~published-and-available blob data together with ordering evidence, and (ii)~accountable hard finality of a header chain. Ordering rests on the \textbf{Bullshark anchor commit rule} rather than heuristic ``frontier'' heuristics. Macro-finality employs \textbf{adaptive subnet BLS aggregation}: flat gossip for smaller validator sets ($N < 500$), subnet-based aggregation for larger sets ($N > 1000$), and interpolated behavior between those thresholds. A mandatory \textbf{two-tier finality model} combines \emph{Fast Execution Finality} ($\approx$5--10\,s) via Casper~FFG-style 2-chain justification on macro checkpoints with \emph{Sovereign Epoch Finality} ($\approx$60\,min) via Bitcoin Taproot checkpointing and six confirmations---the tier required for validator unbonding and for high-value bridge releases.

This whitepaper presents a protocol design. Deploying it in production would require full implementation and independent security review. Performance figures are engineering targets, not guaranteed measurements.

\paragraph{Cryptography.} The stack treated here is \emph{classical}: BLS12-381 for aggregation and ECVRF (Edwards25519, RFC~9381) for private sortition. Any future post-quantum transition should replace primitives in a single coordinated cryptographic migration so that security claims remain aligned across all layers.
\end{abstract}

\tableofcontents
\cleardoublepage

\chapter{Design Foundations}

\section{Architectural commitments}
The following principles shape the whole system:
\begin{enumerate}[leftmargin=*]
  \item \textbf{Honest cryptography labeling:} Leader election uses ECVRF (Edwards25519). Aggregate voting uses BLS12-381. The protocol does not claim post-quantum security while those primitives remain in use; a later upgrade would replace them coherently end-to-end.
  \item \textbf{Throughput and sampling:} Without data availability sampling (DAS), sustained DA throughput is \textbf{capped at 5\,MB/s} so full nodes on typical home links can replicate all data. Raising throughput into the tens of MB/s requires DAS---for example 2D Reed--Solomon with KZG or RLNC-based sampling---and should not be paired with full replication alone.
  \item \textbf{Adaptive subnetting:} Macro vote aggregation scales with active set size: flat gossip at small scale, subnet aggregation at large scale, and leaderless gossip fallback under proposer failure (Chapter~\ref{ch:macro}).
  \item \textbf{Finality boundaries:} Bitcoin checkpointing (Babylon/Pikachu-style) is part of the base design. Fast execution finality and sovereign epoch finality are distinct API-visible states.
  \item \textbf{Multi-layer anti-Sybil:} Beyond a 5\% per-key stake cap, validators declare network placement (ASN, cloud, region), incur concentration-based reward decay, and may participate in DKG-based key-origin accountability as that machinery matures.
\end{enumerate}

\chapter{Executive Summary}

\section{Core design choices}
\begin{enumerate}[leftmargin=*]
  \item \textbf{Narrow scope:} No execution layer; DA plus accountable header-chain finality only.
  \item \textbf{Bullshark anchor rule:} Causal, quorum-based commits replace informal frontier rules.
  \item \textbf{Adaptive macro aggregation:} Flat, subnet, or leaderless modes instead of a fixed subnet count.
  \item \textbf{Permissionless membership:} Activation and exit queues, withdrawal gated on sovereign finality, churn limits, beacon refresh.
  \item \textbf{Explicit API tiers:} \texttt{accepted}, \texttt{soft\_confirmed}, \texttt{finalized}, \texttt{epoch\_finalized}.
  \item \textbf{Quantitative committee analysis:} Concrete capture probabilities for micro committees.
  \item \textbf{Bitcoin anchoring:} Two finality tiers with clear operational meaning.
  \item \textbf{Anti-Sybil layering:} Declarations, decay, and DKG fingerprinting where enabled.
\end{enumerate}

\section{Relationship to established patterns}
Proof-of-stake under partial synchrony, accountable safety under $< \frac{1}{3}$ Byzantine stake, VRF-based private sortition, Casper~FFG-style finality, and light-client sync committees are well-studied ingredients; LUA-DAG composes them with the DAG availability layer and the two-tier finality model above.

\section{Performance targets (indicative)}
Table~\ref{tab:perf-targets} summarizes headline targets. The 5\,MB/s cap under full replication reflects the choice that light clients initially trust the sync committee for DA (no sampling); full nodes carry the full dataset.

\begin{table}[htbp]
  \centering
  \caption{Engineering performance targets (indicative; depend on implementation and deployment).}
  \label{tab:perf-targets}
  \begin{tabular}{@{}lll@{}}
    \toprule
    \textbf{Metric} & \textbf{Full replication (no DAS)} & \textbf{With DAS} \\
    \midrule
    DA throughput (sustained) & 5\,MB/s (hard cap) & 30--100\,MB/s \\
    Soft-confirm latency p95 & 0.5--1.5\,s & 0.5--1.5\,s \\
    Fast execution finality p95 & 5--10\,s & 5--10\,s \\
    Sovereign epoch finality p95 & $\sim$60\,min (6 BTC blocks) & $\sim$60\,min \\
    \bottomrule
  \end{tabular}
\end{table}

\paragraph{Validator hardware (full-replication tier).} A reasonable target is 8 vCPU, 32\,GB RAM, 1\,TB NVMe, and 100\,Mbps connectivity, aligned with the 5\,MB/s cap and consumer-grade uplinks.

\chapter{Positioning, Scope, and Success Metrics}
\label{ch:positioning}

\section{Value proposition}
LUA-DAG offers \textbf{published, available, and finalized blob data} as a service. Customers are rollups, app-chains, and bridges---not end users of decentralized applications on a shared execution layer.

\subsection{Inputs and outputs}
Rollups supply: (i)~blobs (serialized transaction batches), optionally (ii)~state-root commitments or fraud-proof references. The protocol returns:
\begin{itemize}[leftmargin=*]
  \item Availability evidence (DAG vertex certificates).
  \item Ordering evidence (anchor commit within each rollup namespace).
  \item Hard-finality evidence (macro checkpoint with BLS aggregate).
  \item Slashable evidence when safety assumptions are violated.
\end{itemize}

Rollups retain execution, cross-rollup atomicity (if desired), and mempool policy. LUA-DAG \textbf{does not} act as a universal shared sequencer; it commits blob order \emph{within} a namespace.

\section{Explicit non-goals}
\begin{itemize}[leftmargin=*]
  \item Not a generic L1 (no EVM, Move, SVM, etc.).
  \item Not a mandatory shared sequencer.
  \item MEV: default off-chain to rollups; optional fairness mode (Fino-style integration) may be opted in.
  \item Not ``light DA'' in the full-replication configuration: absent sampling, throughput stays capped as above.
  \item Not post-quantum under the classical BLS and ECVRF assumptions stated earlier.
  \item Not a live restaking AVS: native PoS stake grounds safety/liveness; Bitcoin does not vote in MacroQC.
\end{itemize}

\section{Key performance indicators (adversarial testnet)}
High-level testnet success criteria include: fast execution finality p95 $< 10$\,s under nominal WAN; degraded p95 $< 20$\,s with $\frac{1}{4}$ stake offline and 5\% loss; sovereign epoch finality p95 $< 90$\,min; DA throughput $> 5$\,MB/s with 200 validators and mixed blob sizes; soft-confirm p95 $< 2$\,s; weak-subjectivity sync $< 30$\,min on 100\,Mbps home links; mobile header verify $< 5$\,ms; storage growth $< 100$\,GB/validator/month at 5\,MB/s; reliable detection of equivocation and configured data-unavailability scenarios; anti-Sybil false positive rate $< 1\%$ baseline with $\geq 95\%$ true positives on injected six-node splits.

\section{Two-tier finality model}
\label{sec:finality-tiers}
Table~\ref{tab:finality-tiers} defines user-visible tiers. The split exists because accountable slashing loses force once stake has exited: an external PoW anchor closes the game-theoretic gap for validator rotation and large bridge flows.

\begin{table}[htbp]
  \centering
  \caption{Finality tiers.}
  \label{tab:finality-tiers}
  \begin{tabular}{@{}p{3.6cm}p{2.2cm}p{4.2cm}p{4.0cm}@{}}
    \toprule
    \textbf{Tier} & \textbf{Latency} & \textbf{Mechanism} & \textbf{Typical use} \\
    \midrule
    Soft confirmation & 0.5--1.5\,s & MicroQC (Bullshark + micro committee) & UI preview, revertible fees \\
    Fast execution finality & 5--10\,s & MacroQC + Casper 2-chain rule & Rollup tx, capped bridging \\
    Sovereign epoch finality & $\sim$60\,min & BTC Taproot checkpoint + 6 conf & Unbond, rotation, large bridge \\
    \bottomrule
  \end{tabular}
\end{table}

\chapter{System Model and Threat Model}

\section{Validators and stake}
Let $\mathcal{V} = \{v_1,\ldots,v_N\}$ with stake weights $w_i \geq w_{\min}$. Active stake in epoch $e$ is $W_e = \sum_{i \in \mathrm{active}_e} w_i$. Per-validator weight is capped at $w_{\max} = 0.05\, W_e$; excess stake does not increase voting power (it is ``burned'' for weight purposes).


\subsection{Anti-Sybil obligations}
To make the 5\% cap economically meaningful under permissionless joining, validators must:
\begin{enumerate}[leftmargin=*]
  \item \textbf{Declare network identity} at registration: ASN, cloud provider (if any), region; committed on-chain; changes require re-attestation.
  \item \textbf{Submit DKG-fingerprint commitment} when participating in the optional registry; enforceability can strengthen as ceremony tooling matures.
  \item \textbf{Accept correlation-based reward decay} (Section~\ref{sec:incentives-decay}).
\end{enumerate}
Undeclared validators are treated as maximally concentrated for decay purposes.

\section{Network assumptions}
The network is \textbf{partially synchronous}: there exist unknown $\Delta$ and GST such that after GST every message between correct nodes arrives within $\Delta$. Before GST the adversary controls scheduling. Topology is gossip-based with eager push for headers and votes and pull-on-demand for chunks---no rigid relay tree.

\section{Trust assumptions for cryptography}
\begin{itemize}[leftmargin=*]
  \item Hash functions (SHA-256, BLAKE3): collision and second-preimage resistance.
  \item BLS12-381: EUF-CMA under classical assumptions; rogue-key resistance via proof-of-possession.
  \item ECVRF: pseudorandomness and operational bias resistance adequate for stake-weighted sortition (Algorand/Polkadot precedent), \emph{not} PQ-safe.
  \item Bitcoin PoW: standard six-block confirmation rule and economic hardness against deep reorgs.
  \item No synchronized clocks; only local timeouts.
\end{itemize}

\section{Threat summary}
Safety holds for Byzantine stake $f < W/3$; liveness after GST requires honest online stake $> 2W/3$. Adaptive corruption after seeing randomness is allowed subject to epoch-bounded rates. Long-range attacks are mitigated by one-week weak subjectivity plus default Bitcoin checkpoints. Data withholding is addressed by certified vertices and retrieval challenges (Chapter~\ref{ch:da}). Sybil via stake split faces decay and (later) DKG evidence.

\chapter{Architecture Overview}
\label{ch:arch}

\section{Taxonomic placement}
LUA-DAG is an \textbf{ebb-and-flow} style design~\cite{neu2021ebb}: a dynamically available synchronous-feeling layer (DAG availability + micro-ordering) paired with a partially synchronous finality gadget (macro layer), analogous in spirit to RLMD-GHOST + finality gadgets and Ethereum three-slot finality narratives. Layer~4 adds a \textbf{Sovereign Anchor}: PoS-checkpointed-into-PoW, following Babylon~\cite{tas2022babylon} and Pikachu~\cite{azouvi2022pikachu}.

\section{Four layers}
Figure~\ref{fig:layers} summarizes the layering. Each layer exposes a single narrow interface upward (Section~\ref{sec:interfaces}).

\begin{figure}[htbp]
  \centering
  \begin{lstlisting}[
    frame=none,
    basicstyle=\ttfamily\footnotesize,
    columns=fixed,
    keepspaces=true,
    breaklines=false
  ]
                         +-------------------------------+
                         | Rollup / app-chain client     |
                         | (execution + own mempool)     |
                         +---------------+---------------+
                                        | submit blob
                                         v
        +----------------------------------------------------------------+
        | Layer 1: AVAILABILITY DAG (Narwhal-class)                      |
        | - chunks, erasure coding; certified vertices (2f+1)            |
        +----------------------------------------------------------------+
        | Layer 2: MICRO-ORDERING (Bullshark anchor commit)              |
        | - ECVRF anchor; shortcut/slow path; MicroCheckpoint + MicroQC  |
        +----------------------------------------------------------------+
        | Layer 3: MACRO-FINALITY (Casper FFG-class)                     |
        | - MacroCheckpoint; adaptive BLS aggregation; 2-chain rule      |
        +----------------------------------------------------------------+
        | Layer 4: SOVEREIGN ANCHOR (default ON)                         |
        | - BTC Taproot; 6 conf -> epoch_finalized                       |
        +----------------------------------------------------------------+
                                         v
                         +-------------------------------+
                         | Rollup / bridge / light client|
                         +-------------------------------+
  \end{lstlisting}
  \caption{Conceptual four-layer stack (ASCII).}
  \label{fig:layers}
\end{figure}

\section{Inter-layer interfaces}
\label{sec:interfaces}
\begin{itemize}[leftmargin=*]
  \item L1$\to$L2: \texttt{causal\_set(round\_cut)} $\mapsto$ certified vertices up to a round cut.
  \item L2$\to$L3: \texttt{micro\_head()} $\mapsto$ \texttt{MicroCheckpoint} with roots and micro QC.
  \item L3$\to$consumer: \texttt{latest\_finalized()} $\mapsto$ (\texttt{MacroHeader}, \texttt{MacroQC}).
  \item L3+BTC: \texttt{latest\_epoch\_finalized()} adds \texttt{BitcoinAnchorProof}.
  \item Rollup API: \texttt{blob\_status(blob\_id)} in $\{$\texttt{submitted}, \texttt{accepted}, \texttt{ordered}, \texttt{soft\_confirmed}, \texttt{justified}, \texttt{finalized}, \texttt{epoch\_finalized}$\}$.
\end{itemize}

\section{Typical timeline (illustrative)}
Blob submission to soft confirmation can be sub-second to $\sim$1.5\,s under healthy conditions. Macro windows ($W=8$ micro-slots with $\sim$250\,ms rounds) yield $\sim$8\,s windows; two justified generations produce fast execution finality in $\sim$10--16\,s worst case. Bitcoin anchoring typically tracks epoch boundaries on the order of tens of minutes plus six confirmations ($\sim$60\,min sovereign latency).

\section{Blob lifecycle}
States advance monotonically except the deliberate revert window from \texttt{accepted} to \texttt{soft\_confirmed} semantics: soft confirmation may revert under macro stalls if \texttt{lock\_macro} is violated (a protocol bug class) or during slashable macro forks. Section~\ref{sec:cross-layer} details the \texttt{lock\_macro} invariant.

\chapter{Data Structures}
\label{ch:data}

Table~\ref{tab:structs} lists primary structures. Deliberately, no object named ``block'' appears: execution, DA, ordering, and finality are separated.

\begin{table}[htbp]
  \centering
  \caption{Core structs (field lists abbreviated).}
  \label{tab:structs}
  \small
  \begin{tabular}{@{}p{3.2cm}p{10.5cm}@{}}
    \toprule
    \textbf{Struct} & \textbf{Role / main fields} \\
    \midrule
    \texttt{Blob} & Namespace, bytes, Merkle commitment over chunks \\
    \texttt{Chunk} & Fragment after erasure coding + Merkle proof \\
    \texttt{IngressReceipt} & Slashable evidence for \texttt{accepted} \\
    \texttt{Vertex} & Round, author, parents ($\geq 2f{+}1$), blob refs, sig \\
    \texttt{CertifiedVertex} & Vertex + $2f{+}1$ BLS quorum \\
    \texttt{MicroCheckpoint} & Slot, parent macro, anchor, sub-DAG root, ordered blob root, micro QC \\
    \texttt{MicroQC} & Committee bitmap + BLS aggregate on micro head \\
    \texttt{MacroCheckpoint} & Height, parent hash, micro head, DA root, valset root, epoch, proposer \\
    \texttt{MacroQC} & Height, subnet bitmap, aggregates, validator bitmap, BLS, signed stake \\
    \texttt{MacroHeader} & Compressed light-client header + aggregate \\
    \texttt{SyncCommitteeUpdate} & Epoch handoff for light clients \\
    \texttt{BitcoinAnchorProof} & Macro id, BTC txid, Merkle path, confirmations \\
    \texttt{ValidatorIdentity} & ASN, cloud, region attestation \\
    \texttt{DKGCommitment} & Key-origin fingerprint envelope \\
    \texttt{SlashEvidence} & Kind + pairwise conflicting signed material \\
    \bottomrule
  \end{tabular}
\end{table}

\chapter{Layer 1: Availability DAG}
\label{ch:da}

\section{Certified vs.\ uncertified DAG}
\textbf{Baseline deployment:} certified Narwhal/Bullshark-style DAG---parents must already be certified ($2f{+}1$ signatures). This adds roughly one round compared with some uncertified designs but simplifies availability arguments. Uncertified fast paths analogous to Mysticeti remain an optional direction once Adelie-class mitigations are integrated.

\textbf{Shoal-inspired enhancement:} leader reputation and pipelining reduce reliance on timeouts in the common case.

\section{Vertex creation and certification}
Each validator in round $r$ waits for $\geq 2f{+}1$ certified parents from $r-1$, batches blob references subject to \texttt{MAX\_VERTEX\_PAYLOAD} (default 256\,KB here), signs, and gossips. Peers verify signatures, parent certification, and \emph{at least one} chunk availability proof per referenced blob on the fast path; full chunk coverage is enforced via custody and challenges. Quorum signatures promote a \texttt{CertifiedVertex}.

\section{Erasure coding}
Blobs split into $k$ data chunks expanded to $n=2k$ via Reed--Solomon (rate $\frac{1}{2}$). Chunk size 32\,KB baseline; commitment is Merkle root over chunk hashes. Gossip load-balances bytes across validators ($\sim 1/N$ share each).

\subsection{Throughput cap under full replication}
\texttt{THROUGHPUT\_HARD\_CAP} = 5\,MB/s enforced via payload caps and a base-fee targeting $\sim$50\% steady utilization ($\sim$2.5\,MB/s typical). Sustaining tens of MB/s \emph{without} DAS implies multi-TB/day full-node bandwidth and risks centralization.

\subsection{Transition to data availability sampling}
2D Reed--Solomon + KZG commitments (or RLNC-DAS) become appropriate when raising throughput; the vertex schema may add \texttt{kzg\_commitment} with backward compatibility for Merkle-only vertices during migration.

\section{Garbage collection}
\textbf{Hot} rounds (near head): retain full chunks. \textbf{Warm}: keep commitments + at least one chunk. \textbf{Cold}: may delete bytes; archives opt-in.

\section{Custody and retrieval challenges}
\label{sec:challenges}
Custody validators for chunk $(\texttt{blob\_id}, i)$ are
\begin{equation}
  \mathrm{custody}(\texttt{blob\_id}, i) = \bigl\{\, v_j : H(\texttt{pubkey}_j \| \texttt{blob\_id} \| i) \bmod N < K_{\mathrm{custody}} \,\bigr\},
\end{equation}
with default $K_{\mathrm{custody}} = 2f{+}1$. Obligations: store in hot tier, answer challenges within $T_{\mathrm{retrieve}} = 30$\,s, retain Merkle proofs in warm tier.

Challengers broadcast \texttt{Challenge}; silence from a supermajority of custody validators yields \texttt{UnavailabilityEvidence} assembled from signed gossip witness statements (validators need not inspect chunk correctness---only absence of valid gossip responses). Successful proofs slash non-responders (default 5\%) and, if the whole custody set fails, mark the blob unavailable and slash the ingress validator 5\%.

Anti-spam: minimum interval between challenges, fee floor refunded on successful slash.

\chapter{Layer 2: Micro-Ordering (Bullshark)}
\label{ch:micro}

\section{Waves and anchors}
Waves span four rounds: $4w,\ldots,4w{+}3$. Randomness beacon
\begin{equation}
  R_w = H\bigl(R_{w-1} \| \text{MacroQC of latest finalized}\bigr).
\end{equation}
Each validator computes $y_i = \mathrm{ECVRF}_i(R_w \| \text{``anchor''})$; the anchor minimizes $y_i \cdot W / (w_i \cdot \mathrm{rep}_i)$ with Shoal-style reputation $\mathrm{rep}_i \in [0.8, 1.2]$, narrowing the bias surface from wider reputation windows. Identity remains hidden until the anchor publishes at round $4w$.

\section{Commit rules}
\textbf{Shortcut (two rounds):} Anchor $A_w$ at $4w$ commits if $\geq 2f{+}1$ certified vertices at $4w{+}1$ include $A_w$ as parent.

\textbf{Slow path (four rounds):} If shortcut fails, validators vote at $4w{+}2$; at $4w{+}3$, $2f{+}1$ support commits, or skip advances the wave.

\section{Linearization}
Given committed anchor $A_w$, $\mathrm{Closure}(A_w)$ excludes vertices in the previous anchor's closure. Deterministic topological sort (BFS from $A_w$, tie-break by vertex hash) orders vertices; blob refs group by namespace preserving local order. Roots \texttt{committed\_sub\_dag\_root} and \texttt{ordered\_blob\_refs\_root} summarize the wave.

\section{MicroQC and soft confirmation}
A micro committee (default size 256, VRF-sampled) signs the \texttt{MicroCheckpoint} hash; aggregation at $\lceil \frac{2}{3} C_{\mathrm{micro}} \rceil$ forms \texttt{MicroQC}. Presence of \texttt{MicroQC} marks blobs \texttt{soft\_confirmed} with API flag \texttt{revert\_risk = true}.

\section{Why Bullshark suffices here}
Asynchronous networks lack a global DAG snapshot; Bullshark commits via causal quorum linkage propagating with messages. Correct nodes that see supporting vertices must eventually see the anchor and its ancestors, forcing agreement on commits without wall-clock synchronicity.

\chapter{Layer 3: Macro-Finality}
\label{ch:macro}

\section{Cadence and proposer}
Every $W=8$ micro-slots, emit a \texttt{MacroCheckpoint}. Macro proposer selection mirrors micro anchors with beacon $R^{\mathrm{macro}}_h = H(R^{\mathrm{macro}}_{h-1} \| \mathrm{MacroQC}_{h-1})$, primary and backup proposers, timeout $T_{\mathrm{macropropose}} = 4$\,s.

\section{Adaptive aggregation (Modes 0, A, B)}
\label{sec:modes}
Let $N_e$ be active validators in epoch $e$. Define
\begin{equation}
  K_e =
  \begin{cases}
    0, & N_e < 500 \quad \text{(Mode 0: flat gossip)}\\
    \bigl\lceil N_e / 128 \bigr\rceil, & 500 \leq N_e \leq 1000 \quad \text{(interpolated)}\\
    \bigl\lceil N_e / 128 \bigr\rceil, & N_e > 1000 \quad \text{(Mode A: subnets)}
  \end{cases}
\end{equation}
capped at $K_{\max} = 32$. Subnet membership rotates per-epoch via
\begin{equation}
  \mathrm{subnet}(v_i,e) = H(\texttt{pubkey}_i \| R^{\mathrm{macro}}_{\mathrm{epoch\_start}(e)}) \bmod K_e.
\end{equation}

\paragraph{Mode 0 ($K_e=0$).} Validators sign \texttt{MacroProposal} and gossip partials; proposer aggregates until $\geq 2W/3$ stake or timeout. Failure triggers Mode~B.

\paragraph{Mode A ($K_e \geq 4$).} Subnet aggregators collect partials, publishing subnet aggregates at internal $2/3$ stake \emph{or} at deadline $T_{\mathrm{subnet}} = T_{\mathrm{macropropose}}/2$. Macro proposer unions subnets whose verified aggregates sum to $\geq 2W/3$ stake.

\paragraph{Mode B (fallback).} Validators gossip partials; any node may form a candidate \texttt{MacroQC}. Canonical winner maximizes signed stake; ties break by lexicographic validator (or subnet) bitmap. Re-broadcast rules yield $O(\log N)$ gossip rounds convergence.

\section{Two-chain finality}
$C_h$ is \textbf{justified} with a valid \texttt{MacroQC}. It is \textbf{finalized} (fast execution tier) when some $C_{h+1}$ is also justified and $\texttt{parent\_height\_hash}(C_{h+1}) = H(C_h)$. Expected latency $\approx 2W \times (\text{micro slot duration})$.

\section{Slashing (macro layer)}
Equivocating macro votes, surrounding votes, double propose: 100\% slash. Micro double-votes: 50\%. Proven data unavailability: 5\% per incident (annual soft-cap 50\%). Inactivity leak: if no finality for four macro windows, offline/mis-voting stake bleeds $\approx$0.5\% per window until supermajority restores progress.

\chapter{Light Clients, Checkpoints, and Bitcoin Anchoring}

\section{Sync committee}
Each epoch (e.g.\ 1024 macro heights) samples \texttt{SYNC\_COMMITTEE\_SIZE} = 512 slots with replacement, weighted by stake. The committee signs every \texttt{MacroHeader}; verification is one pairing + bitmap parse ($<5$\,ms mobile target). The committee \emph{does not} finalize---it is a gateway for resource-constrained verifiers.

\section{Weak subjectivity}
\texttt{WEAK\_SUBJECTIVITY\_PERIOD} = one week. Checkpoints must be fresh from trusted sources ($\geq 3$ recommended). Bitcoin anchoring narrows long-range exposure for validating nodes; weak subjectivity still bootstraps initial sync.

\section{Bitcoin checkpointing (default ON)}
Vigilante relayers aggregate \texttt{BitcoinCheckpoint} = $(\text{epoch}, \text{macro\_height}, \text{macro\_hash}, \text{validator\_set\_root})$, collect $\geq 2W/3$ BLS signatures, embed 64-byte payload in Taproot script-path spend, broadcast with adaptive fees, and publish \texttt{BitcoinAnchorProof} after six confirmations---promoting matching macro heights to \texttt{epoch\_finalized}.

\paragraph{Economics and failure modes.} Relayers earn \texttt{BTC\_RELAY\_REWARD} (default 5\% of macro proposer reward class). If all relayers stall, rollup fast finality continues but validator exits halt. Deep BTC reorgs $>6$ invalidate proofs until re-anchoring (explicit accountable-halt narrative).

\section{Bridge value cap}
\texttt{BRIDGE\_VALUE\_CAP} defaults to $1000 \times \texttt{MIN\_STAKE}$; bridges may go lower. Withdrawals above the cap require \texttt{epoch\_finalized}.

\chapter{Permissionless Membership}

\section{Activation and churn}
Validators deposit, queue under \texttt{MAX\_ACTIVATION\_PER\_EPOCH} (default 4). Exits queue at \texttt{MAX\_EXIT\_PER\_EPOCH}. While exiting after acceptance in macro state, validators stop voting but remain slashable until sovereign finality.

\section{Withdrawal delay}
Stake unlocks only after the containing macro checkpoint reaches \texttt{epoch\_finalized} with \texttt{WITHDRAWAL\_DELAY} = six Bitcoin confirmations ($\sim$60\,minutes), replacing multi-week PoS-only delays when anchoring is healthy.

\section{Randomness beacon refresh}
MacroQC-chained beacons resist grinding below $1/3$ stake. Inclusion incentives and subnet pre-publication prevent proposers from hiding visible aggregates.

\section{Anti-Sybil mechanisms}
\label{sec:antisybil}
Validators submit \texttt{ValidatorIdentity}. Concentration score:
\begin{equation}
  \mathrm{concentration\_score}(v_i) =
    \alpha\, s_{\mathrm{ASN}}(v_i) + \beta\, s_{\mathrm{cloud}}(v_i) + \gamma\, s_{\mathrm{voting}}(v_i),
\end{equation}
defaults $(\alpha,\beta,\gamma) = (0.4, 0.4, 0.2)$. Declaring \texttt{self\_hosted} zeros cloud concentration term.

DKG registry participation may begin as opt-in infrastructure and can become mandatory for new activations under governance as ceremony protocols stabilize. The slash escalator is $\texttt{DKG\_SLASH\_BASE} \times 2^{n-1}$ for $n$ colluding origins.

\chapter{Incentives and Accountability}
\label{sec:incentives-decay}

\section{Reward split (illustrative parameters)}
Base 28\%, vertex authoring 15\%, anchor proposing 10\%, micro committee 5\%, macro voting 28\%, macro proposing 5\%, aggregation 4\%, Bitcoin relay 5\%---tunable by governance.

\section{Fee market}
Ingress validators capture the bulk of fees; macro proposer share incentivizes DA root inclusion; 20\% burn. EIP-1559-style base fee tracks target utilization.

\section{Anti-Sybil decay on rewards}
\begin{equation}
  R_{\mathrm{actual}}(v_i) =
    R_{\mathrm{base}}(v_i) \times \bigl(1 - \texttt{REWARD\_DECAY\_RATE} \times \mathrm{concentration\_score}(v_i)\bigr),
\end{equation}
with \texttt{REWARD\_DECAY\_RATE} default 1. Bootstrap grace halves decay for 30 epochs when $<50$ validators.

\chapter{Security Analysis}

\section{Macro safety (sketch)}
If $f < W/3$, two finalized conflicting checkpoints imply two $\geq 2W/3$ vote sets intersecting in $\geq W/3$ stake, yielding Casper-class slashable surround/conflict evidence.

\section{Liveness (sketch)}
After GST, honest anchors appear with sufficient probability; shortcut commits progress waves; honest macro proposers assemble QCs; consecutive justified heights finalize. Termination is eventual consistent with partial synchrony and FLP.

\section{Micro committee capture probability}
For committee size $C$ and Byzantine stake fraction $\beta$,
\begin{equation}
  P(\text{capture}) = \sum_{k=\lceil C/3\rceil}^{C} \binom{C}{k} \beta^k (1-\beta)^{C-k}.
\end{equation}
Near $\beta \approx 1/3$, finite committees never make soft confirmation as strong as macro finality---hence API separation.

\section{Attack matrix (selected rows)}
Long-range: weak subjectivity window + BTC anchor. Data withholding: challenges + slashing. Anchor DoS: slow path + reputation. Macro proposer DoS: Mode~B. Subnet capture: VRF partitions + quorum across surviving subnets. Micro-flush: \texttt{lock\_macro} (Section~\ref{sec:cross-layer}). Sync committee lying: no finality power; cross-check + BTC. Post-quantum: outside the classical assumptions of this document.

\section{Cross-layer interaction: \texttt{lock\_macro}}
\label{sec:cross-layer}
Micro anchors must build on DAG views whose macro parent equals the latest locally justified macro checkpoint \emph{without downgrade}. If a peer advertises a higher \texttt{lock\_macro}, verify its \texttt{MacroQC} before certifying. Lower \texttt{lock\_macro} anchors are rejected. Micro checkpoints record \texttt{parent\_macro} explicitly. Macro forks freeze soft-finality semantics until slashable resolution.

\chapter{Performance Plan}
Parameter sweeps span validator counts 50--500, WAN topologies, round 100--500\,ms, macro windows 4--16, committee sizes, subnet counts, blob mixes, erasure rates, and Byzantine fractions. Measure throughput, latency percentiles, bandwidth, CPU, storage, recovery times, fairness, and security signal rates.

Expected bottlenecks: chunk gossip bandwidth, BLS verification volume, disk IOPS, subnet aggregation latency at scale.

\chapter{Roadmap, Risks, and Open Questions}

\section{Phased delivery}
A realistic engineering program proceeds from formal modeling (for example TLA+) through layered prototypes: availability DAG, Bullshark micro-ordering, macro finality with adaptive aggregation, Bitcoin anchoring and light clients, permissionless membership with anti-Sybil instrumentation, adversarial testnets, and external audit. Bitcoin checkpoint and anti-Sybil workstreams add parallel calendar time relative to a consensus-only core.

\section{Risk register}
Cross-layer bugs, aggregation correctness, Bitcoin operations, sync committee deception, anti-Sybil tuning, and adoption competition head the list.

\section{Future extensions}
Mandatory DAS and mandatory DKG for new activations are natural hardening steps once tooling is ready; an optional uncertified DAG mode may follow the same maturity path. A coordinated cryptographic upgrade could move aggregation and leader election to post-quantum primitives. Explicitly out of scope in this document: economic security borrowed from restaking as the live liveness backbone, universal shared sequencing for all rollups, and native IBC-style transport.

\chapter{Comparison with Related Systems}
Relative to modular DA networks, LUA-DAG emphasizes DAG-shaped availability with load spreading, explicit two-tier finality, and accountable macro votes backed by a sovereign Bitcoin anchor. Compared with Celestia: DAG load balancing and the sovereign epoch tier. Compared with Avail: different timing on sampling versus full replication, plus macro finality and BTC anchoring choices here. Compared with Mysticeti/Sui-class systems: DA-only scope rather than shared execution. Compared with Acki Nacki: full-validator macro accountability rather than probabilistic per-block committees. Compared with Babylon: integrated DAG DA layer alongside checkpointing economics.

\chapter{Parameter Reference}
Table~\ref{tab:params} lists representative constants; governance may tune within operational ranges.

\begin{table}[htbp]
  \centering
  \caption{Selected protocol parameters (illustrative defaults).}
  \label{tab:params}
  \small
  \begin{tabular}{@{}l l@{}}
    \toprule
    \textbf{Parameter} & \textbf{Default} \\
    \midrule
    \texttt{ROUND\_DURATION} & 250\,ms \\
    \texttt{MACRO\_WINDOW\_W} & 8 micro-slots \\
    \texttt{MICRO\_COMMITTEE\_SIZE} & 256 \\
    \texttt{SUBNET\_FLAT\_THRESHOLD} / \texttt{FULL} & 500 / 1000 validators \\
    \texttt{MAX\_VERTEX\_PAYLOAD} & 256\,KB \\
    \texttt{MAX\_BLOB\_SIZE} & 1\,MB \\
    \texttt{THROUGHPUT\_HARD\_CAP} & 5\,MB/s \\
    \texttt{GC\_HOT\_HORIZON} / \texttt{WARM} & 200 / 10{,}000 rounds \\
    \texttt{WEAK\_SUBJECTIVITY\_PERIOD} & 7 days \\
    \texttt{WITHDRAWAL\_DELAY} & 6 BTC confirmations \\
    \texttt{T\_MACROPROPOSE} / \texttt{T\_SUBNET} / \texttt{T\_CANONICALIZE} & 4\,s / 2\,s / 8\,s \\
    \texttt{SYNC\_COMMITTEE\_SIZE} & 512 \\
    \texttt{BTC\_CONFIRMATIONS\_FOR\_FINAL} & 6 \\
    \texttt{IDENTITY\_REATTEST\_PERIOD} & 4 epochs \\
    \texttt{REWARD\_DECAY\_RATE} & 1.0 \\
    \texttt{DKG\_SLASH\_BASE} & 20\% \\
    \bottomrule
  \end{tabular}
\end{table}

\chapter*{Glossary (abridged)}
\addcontentsline{toc}{chapter}{Glossary (abridged)}

\begin{description}[leftmargin=0pt,labelindent=0pt,itemsep=0.35em]
  \item[Anchor.] VRF-elected commit pivot for a Bullshark wave.
  \item[Bitcoin anchor proof.] Evidence bundling macro hash into BTC with confirmations.
  \item[Certified vertex.] Vertex certified by $2f{+}1$ signatures---only such vertices may parent new work.
  \item[Fast execution finality.] Casper 2-chain justified macro state ($\sim$seconds-to-ten seconds).
  \item[Sovereign epoch finality.] BTC-confirmed checkpoint state ($\sim$one hour).
  \item[Mode 0 / A / B.] Flat, subnet, or leaderless macro aggregation pathways.
\end{description}

\begin{thebibliography}{99}
\bibitem{neu2021ebb}
Neu, Tas, and Tse,
\newblock \emph{Ebb-and-Flow Protocols: A Resolution of the Availability-Finality Dilemma},
\newblock IEEE S\&P, 2021.

\bibitem{tas2022babylon}
Tas \emph{et al.},
\newblock \emph{Bitcoin-Enhanced Proof-of-Stake Security: Possibilities and Impossibilities (Babylon)},
\newblock IEEE S\&P, 2023.

\bibitem{azouvi2022pikachu}
Azouvi \emph{et al.},
\newblock \emph{Pikachu: Securing PoS Blockchains from Long-Range Attacks by Checkpointing into Bitcoin PoW},
\newblock 2022.

\bibitem{danezis2022narwhal}
Danezis \emph{et al.},
\newblock \emph{Narwhal and Tusk: A DAG-based Mempool and Efficient BFT Consensus},
\newblock EuroSys, 2022.

\bibitem{spiegelman2022bullshark}
Spiegelman \emph{et al.},
\newblock \emph{Bullshark: DAG BFT Protocols Made Practical},
\newblock CCS, 2022.

\bibitem{buterin2017casper}
Buterin and Griffith,
\newblock \emph{Casper the Friendly Finality Gadget},
\newblock arXiv:1710.09437, 2017.

\bibitem{dwork1988partial}
Dwork, Lynch, Stockmeyer,
\newblock \emph{Consensus in the Presence of Partial Synchrony},
\newblock JACM, 1988.

\bibitem{fischer1985flp}
Fischer, Lynch, Paterson,
\newblock \emph{Impossibility of Distributed Consensus with One Faulty Process},
\newblock JACM, 1985.
\end{thebibliography}

\appendix

\chapter{Blob State Machine (formal summary)}
States: $\texttt{submitted} \to \texttt{accepted} \to \texttt{ordered} \to \texttt{soft\_confirmed} \to \texttt{justified} \to \texttt{finalized} \to \texttt{epoch\_finalized}$. Typical triggers: ingress receipt; inclusion in a certified vertex; membership in an anchor's causal closure; justification by a macro checkpoint containing the relevant micro head; the two-chain finality rule; Bitcoin anchor proof with sufficient confirmations. Revert risks: \texttt{soft\_confirmed} and \texttt{justified} may regress during slashable macro forks; \texttt{finalized} is reversible only through accountable slashing inside the unbond window; \texttt{epoch\_finalized} would require a deep Bitcoin reorg to invalidate.

\chapter{Committee Capture Reference Table}
Numeric values for committee capture appear in Appendix~\ref{app:capture}.

\chapter{Bibliographic notes}
The bibliography lists primary references. BFT, DAG-based consensus, data availability sampling, finality gadgets, and PoS long-range security each have large literatures; readers should follow citations from the core papers below into those research lines.

\chapter{Adversarial Testnet KPIs}
\label{app:kpi}

\begin{table}[htbp]
  \centering
  \caption{Key performance indicators for adversarial testnets under the full-replication throughput cap.}
  \small
  \begin{tabular}{@{}p{7.8cm}p{6.8cm}@{}}
    \toprule
    \textbf{KPI} & \textbf{Pass threshold} \\
    \midrule
    Fast execution finality latency p95 & $< 10$\,s (200 validators, 5 WAN regions, 0 Byzantine) \\
    Fast execution finality p95 (degraded) & $< 20$\,s ($\frac{1}{4}$ stake offline, 5\% packet loss) \\
    Sovereign epoch finality p95 & $< 90$\,min (6 BTC confirmations + propagation) \\
    DA throughput sustained & $> 5$\,MB/s (200 validators; 64\,KB--1\,MB blob mix) \\
    Soft-confirm latency p95 & $< 2$\,s \\
    Weak-subjectivity state sync & $< 30$\,min on 100\,Mbps residential uplink \\
    Light client header verify & $< 5$\,ms (Snapdragon~8-class mobile) \\
    Storage growth & $< 100$\,GB/validator/month at 5\,MB/s sustained \\
    Slashable evidence (equivocation) & 100\% detection \\
    Slashable evidence (DA test scenario) & $> 99\%$ \\
    Anti-Sybil alarms & $< 1\%$ false positives; $\geq 95\%$ true positives (6-node split inject) \\
    \bottomrule
  \end{tabular}
\end{table}

\noindent\textbf{Out of scope for these KPIs:} application-level TPS; optional MEV-fairness modes; cross-rollup atomic latency; sustained DA throughput above 5\,MB/s without introducing sampling.

\chapter{Reward Pool Split (Default Parameters)}
\label{app:rewards}

\begin{table}[htbp]
  \centering
  \caption{Epoch reward decomposition (governance-tunable).}
  \begin{tabular}{@{}l r p{7.5cm}@{}}
    \toprule
    \textbf{Reward type} & \textbf{\%} & \textbf{Condition (summary)} \\
    \midrule
    Base reward & 28 & Online; signatures in $\geq$80\% of certified vertices in epoch \\
    Vertex authoring & 15 & Vertices certified on time \\
    Anchor proposing & 10 & Selected as micro anchor and wave commits \\
    Micro committee & 5 & Sampled into micro committee; votes MicroQC \\
    Macro voting & 28 & Votes MacroQC on $\geq$95\% of macro heights \\
    Macro proposing & 5 & Macro proposals become justified \\
    Subnet / flat aggregation & 4 & Timely subnet or Mode~0 aggregation \\
    Bitcoin vigilante relay & 5 & Relayer; checkpoint reaches required confirmations \\
    \bottomrule
  \end{tabular}
\end{table}

\chapter{Attack Matrix and Mitigations (Summary)}
\label{app:attacks}

\begin{longtable}{@{}p{3.6cm} p{4.6cm} p{6.0cm}@{}}
  \toprule
  \textbf{Attack} & \textbf{Effect} & \textbf{Mitigation} \\
  \midrule
  \endfirsthead
  \multicolumn{3}{c}{{\tablename\ \thetable{} --- continued}} \\
  \toprule
  \textbf{Attack} & \textbf{Effect} & \textbf{Mitigation} \\
  \midrule
  \endhead
  Long-range PoS fork & Fake history for new syncers & 1-week WS + default BTC checkpoint \\
  Data withholding & Unavailable blob under soft confirm & Certified vertices; challenges; slashing \\
  Anchor DoS & No shortcut commit & Private sortition; 4-round slow path \\
  Macro proposer DoS & Macro window loss & Mode~B leaderless gossip; timeout \\
  Equivocation & Conflicting checkpoints & 100\% slash; Casper 2-chain \\
  Subnet capture & Bad subnet aggregates & VRF partitions; quorum across subnets \\
  Mode B canonical race & Divergent MacroQC candidates & Stake priority + lex bitmap tie-break \\
  Storage spam & Meaningless DA bytes & Fees; caps; 5\,MB/s full-replication bound \\
  Sync committee lying & Bad headers for light clients & No finality power; full-node checks; BTC anchor \\
  Sybil stake split & Fake decentralization & 5\% cap + decay + DKG path \\
  Unbond + long-range & Slash-evading fork & Stake locked until \texttt{epoch\_finalized} \\
  Bitcoin reorg & Invalid anchor proof & Re-checkpoint; documented revert case \\
  Post-quantum break & Forged classical signatures & Addressed only under a future coordinated crypto migration \\
  \bottomrule
\end{longtable}

\chapter{Micro Committee Capture Probabilities}
\label{app:capture}

Probability that a committee of size $C$ contains at least $\lceil C/3 \rceil$ Byzantine weight when each validator is independently Byzantine with probability~$\beta$:
\[
  P(\text{capture}) = \sum_{k=\lceil C/3 \rceil}^{C} \binom{C}{k} \beta^k (1-\beta)^{C-k}.
\]

\begin{table}[htbp]
  \centering
  \caption{Reference capture probabilities (stake-weighted sampling approximated by Bernoulli $\beta$ per slot).}
  \begin{tabular}{@{}l c c c@{}}
    \toprule
    $C$ & $\beta=0.20$ & $\beta=0.30$ & $\beta=0.33$ \\
    \midrule
    64 & $1.5 \times 10^{-3}$ & 0.18 & 0.42 \\
    128 & $2.7 \times 10^{-5}$ & 0.06 & 0.40 \\
    256 & $1.0 \times 10^{-8}$ & 0.005 & 0.37 \\
    512 & $1.5 \times 10^{-15}$ & $4 \times 10^{-5}$ & 0.34 \\
    1024 & $< 10^{-30}$ & $3 \times 10^{-9}$ & 0.30 \\
    \bottomrule
  \end{tabular}
\end{table}

\noindent With default $C_{\mathrm{micro}}=256$ and $\beta=0.2$, soft confirmation has $\sim 10^{-8}$ committee failure rate---acceptable for UX, not for settlement-grade guarantees (\texttt{finalized} required).

\chapter{Finality Requirements by Use Case}
\label{app:usecases}

\begin{table}[htbp]
  \centering
  \small
  \caption{Mandatory finality tier per integration scenario.}
  \begin{tabular}{@{}p{5.2cm} p{2.4cm} p{5.8cm}@{}}
    \toprule
          \textbf{Use case} & \textbf{Latency} & \textbf{Tier} \\
    \midrule
    UI / fee preview & $<2$\,s & \texttt{soft\_confirmed} \\
    Rollup transaction commit & 5--10\,s & \texttt{finalized} \\
    Cross-rollup messaging & 5--10\,s & \texttt{finalized} \\
    Bridge withdrawal $\leq$ \texttt{BRIDGE\_VALUE\_CAP} & 5--10\,s & \texttt{finalized} \\
    Bridge withdrawal $>$ cap & $\sim$60\,min & \texttt{epoch\_finalized} \\
    Validator unbond / rotation & $\sim$60\,min & \texttt{epoch\_finalized} \\
    Treasury / slashing distribution (long-range sensitive) & $\sim$60\,min & \texttt{epoch\_finalized} \\
    \bottomrule
  \end{tabular}
\end{table}

\end{document}
